\documentclass[12pt]{article}
\usepackage[utf8]{inputenc}
\usepackage[T1]{fontenc}
\usepackage{hyperref}
\title{Tingui-Boto Language Resources for K-12 Curriculum Development}
\author{Compiled by Claude Sonnet 5 (claude-sonnet-5) for Andr\'as Kornai}
\date{}
\begin{document}
\maketitle

\section{Introduction}
Tingui-Boto (ISO 639-3 tgv) belongs to the Tingui-Bot\'o people of Alagoas, Brazil. Glottolog describes the language entry as retired or ``Bookkeeping''. The community speaks Portuguese, and keeps part of an ancestral language, identified as Dzubukuá, for ritual use. Joshua Project calls it a Tikuna dialect and places it in Amazonia, which conflicts with every other source and is not used here. No teaching materials for the language were found.

\section{Primary Educational Resources}
None were found.

\section{Linguistic and Reference Materials}
\begin{itemize}
\item Glottolog entry for ISO code tgv: \url{https://glottolog.org/resource/languoid/iso/tgv}
\item Tingui-Bot\'o people, Wikipedia: \url{https://en.wikipedia.org/wiki/Tingui-Bot%C3%B3_people}
\item Povos Ind\'igenas no Brasil (ISA) page on the Tingui Bot\'o: \url{https://pib.socioambiental.org/en/Povo:Tingui_Bot%C3%B3}
\item A thesis or e-book from the Universidade Federal de Sergipe repository came up in the search. I could not verify its content or relevance, so treat it as a lead only: \url{https://ri.ufs.br/bitstream/riufs/13958/2/FERNANDA_ANDRADE_SILVA_ebook.pdf}
\end{itemize}

\section{Religious and Community Resources}
FUNAI document on the Tingui-Bot\'o Indigenous Land (Olhos d'\'agua do Meio, Feira Grande and Campo Grande, Alagoas): \url{https://www.gov.br/funai/pt-br/arquivos/conteudo/cggam/pdf/2019/ti-tingui-boto-versao-final.pdf}

\section{Digital and Audio Resources}
None were found.

\section{Curriculum Development Context}
Nothing in the sources supports a language curriculum. Any heritage programme would have to work from the community's own ritual vocabulary, and the sources above do not document it. This dossier is deliberately short.
\end{document}